\chapter{The Lebesgue Integral}

\begin{center}\itshape
To add up a pile of coins, you can take them one by one in the order they lie on the table, or sort them first by value
and count each kind. Riemann does the first, Lebesgue the second. Why does sorting by value make limits so much
easier to handle?
\end{center}


Henri Lebesgue liked to explain his integral with a story about paying a debt. (The story is retold everywhere, in several wordings;
we could not trace an exact wording to a document of Lebesgue's, so we paraphrase it.) To pay a sum, you can take the bills and coins out of your
pocket one by one, in the order they come, and hand them over until the total is reached: that is Riemann's integral. Or you can first
empty your pocket, sort the money into heaps of equal value, and pay heap by heap: that is Lebesgue's.

The joke hides a revolution. Bernhard Riemann (1854) cut the axis of the variable into small intervals and added $f(x_i)\,\Delta x_i$.
Lebesgue cut the axis of \emph{values}: the level sets $\{y_j\le f<y_{j+1}\}$ can be any measurable sets, wildly scattered, and he
adds $y_j\,\mu\{y_j\le f<y_{j+1}\}$ (Figure~\ref{fig:sums}). Two things follow. The integral makes sense for any measure: lengths,
probabilities, counting, areas of surfaces. And limits pass under the integral sign under astonishingly mild conditions. Within
four years, Beppo Levi (monotone convergence, 1906), Pierre Fatou (his lemma, 1906) and Lebesgue himself (dominated convergence)
had found the three theorems that every analyst and probabilist now uses daily. Riemann's integral, by comparison, makes you
check uniform convergence, and still cannot integrate the indicator of the rationals.

This chapter keeps three promises. The integral of a continuous function on a segment exists, as Volume~1 admitted, and Lebesgue's
criterion says exactly which bounded functions have a Riemann integral. Uniform integrability, announced in Volume~4,
Chapter~11, is the exact condition for passing to the limit in $L^1$ (Vitali's theorem). And Feynman's trick of Volume~5,
differentiating under the integral sign, becomes a theorem with checkable hypotheses.

\begin{figure}[ht]
\centering
\fig{ch07-sums}
\caption{Riemann sums (left) cut the horizontal axis into intervals and use rectangles. Lebesgue sums (right) cut the vertical axis into levels $y_j$ and use
the measure of the set where $f\ge y_j$, a union of intervals here, any measurable set in general: the integral is $\sum_j(y_{j+1}-y_j)\,\mu\{f\ge y_{j+1}\}$ in the limit.}
\label{fig:sums}
\end{figure}

\section{Measurable functions}

Throughout, $(E,\mathcal A,\mu)$ is a measure space. Functions may take the values $\pm\infty$ when this is convenient; $[0,+\infty]$ carries the obvious order and
topology.

\begin{definition}
A function $f:E\to[-\infty,+\infty]$ is \emph{measurable} if $\{f\le t\}\in\mathcal A$ for every real $t$. A \emph{simple} function is a finite linear combination $s=\sum_ia_i\mathbf 1_{A_i}$ of
indicators of sets $A_i\in\mathcal A$.
\end{definition}

As noted in Chapter~7, $f$ is then measurable for the full Borel $\sigma$-algebra ($\{f\in B\}\in\mathcal A$ for Borel $B$), and continuous functions of measurable functions are
measurable. Indeed the class of $B$ with $\{f\in B\}\in\mathcal A$ is a $\sigma$-algebra containing the half-lines, and if $g:\R^n\to\R$ is continuous, $\{g(f_1,\dots,f_n)<t\}$ is the
preimage of an open set of $\R^n$, a countable union of boxes, each of which pulls back to a finite intersection of sets $\{a_i<f_i\le b_i\}$.

\begin{proposition}\label{prop:measurablefns}
\begin{enumerate}[label=(\alph*)]
\item Sums, products, maxima and minima of finitely many real measurable functions are measurable, and so are $f^+=\max(f,0)$, $f^-=\max(-f,0)$ and $|f|$.
\item If $f_n$ are measurable, so are $\sup_nf_n$, $\inf_nf_n$, $\limsup_nf_n$, $\liminf_nf_n$, and $\lim_nf_n$ where it exists.
\item If $f\ge0$ is measurable, the simple functions $s_n=\min\bigl(n,2^{-n}\lfloor2^nf\rfloor\bigr)$ are measurable, satisfy $0\le s_1\le s_2\le\cdots$, and $s_n\to f$ pointwise; uniformly if $f$
is bounded.
\end{enumerate}
\end{proposition}
\begin{proof}
(a) Continuous functions of measurable functions. (b) $\{\sup_nf_n\le t\}=\bigcap_n\{f_n\le t\}$; $\inf_nf_n=-\sup_n(-f_n)$; $\limsup_n=\inf_m\sup_{n\ge m}$; the set where the limit exists is
$\{\limsup=\liminf\}$, measurable. (c) $s_n$ takes the values $k2^{-n}$, $0\le k\le n2^n$, on the measurable sets $\{k2^{-n}\le f<(k+1)2^{-n}\}$ (and $n$ on $\{f\ge n\}$). Passing from $n$ to $n+1$
halves the grid and raises the cap, so $s_n\le s_{n+1}$. Where $f(x)<n$, $0\le f(x)-s_n(x)<2^{-n}$; where $f(x)=+\infty$, $s_n(x)=n$.
\end{proof}

\section{The integral and the three convergence theorems}

\begin{definition}
For a simple $s=\sum_ia_i\mathbf 1_{A_i}\ge0$ with disjoint $A_i$, let $\int s\,d\mu=\sum_ia_i\mu(A_i)$ (with $0\cdot\infty=0$). For a measurable $f\ge0$, let
\[\int f\,d\mu=\sup\Bigl\{\int s\,d\mu:\ s\text{ simple},\ 0\le s\le f\Bigr\}\in[0,+\infty].\]
For $A\in\mathcal A$, $\int_Af\,d\mu=\int f\mathbf 1_A\,d\mu$.
\end{definition}

The integral of a simple function does not depend on how it is written: two representations with disjoint sets have a common refinement (the
intersections $A_i\cap B_j$), on which both sums agree. It is linear and monotone on simple functions (compute on a common refinement), so the two
definitions agree for simple $f$, and $f\le g$ implies $\int f\le\int g$.

Here is the first of the three theorems, and the engine of the other two. If the functions only go up, nothing can go wrong:
the integral of the limit is the limit of the integrals, even when both are infinite.

\begin{theorem}[monotone convergence]\label{thm:mct6}
If $0\le f_1\le f_2\le\cdots$ are measurable and $f_n\to f$ pointwise, then $\int f_n\,d\mu\to\int f\,d\mu$.
\end{theorem}
\begin{proof}
$f$ is measurable, and $\int f_n$ increases to a limit $L\le\int f$. Conversely let $s=\sum_ia_i\mathbf 1_{A_i}\le f$ be simple and $c\in\,]0,1[$. The sets $E_n=\{f_n\ge cs\}$ increase to $E$: at a point
where $s(x)>0$, $f_n(x)\to f(x)\ge s(x)>cs(x)$, and where $s(x)=0$ every $E_n$ contains $x$. So $\int f_n\ge\int cs\mathbf 1_{E_n}=c\sum_ia_i\mu(A_i\cap E_n)\to c\int s$ by continuity from below.
Hence $L\ge c\int s$; let $c\to1$ and take the supremum over $s$.
\end{proof}

\begin{corollary}\label{cor:linearity6}
For measurable $f,g\ge0$ and $a,b\ge0$, $\int(af+bg)=a\int f+b\int g$. For measurable $f_n\ge0$, $\int\sum_nf_n=\sum_n\int f_n$. If $f\ge0$, then $\nu(A)=\int_Af\,d\mu$ is a measure.
\end{corollary}
\begin{proof}
Approximate $f$ and $g$ by increasing simple functions (Proposition~\ref{prop:measurablefns}(c)); linearity holds for simple functions and passes to the limit by
monotone convergence. The series is the increasing limit of its partial sums. For $\nu$, apply the series statement to $f\mathbf 1_{A_k}$ for disjoint $A_k$.
\end{proof}

\begin{proposition}\label{prop:markov}
Let $f\ge0$ be measurable. (a) (Markov) $\mu(f\ge t)\le\frac1t\int f\,d\mu$ for $t>0$. (b) $\int f=0$ iff $f=0$ a.e. (c) If $\int f<\infty$, then $f<\infty$ a.e. (d) If $f=g$ a.e., then
$\int f=\int g$.
\end{proposition}
\begin{proof}
(a) $t\mathbf 1_{\{f\ge t\}}\le f$. (b) If $\int f=0$, (a) gives $\mu(f\ge\frac1k)=0$ for all $k$, and $\{f>0\}=\bigcup_k\{f\ge\frac1k\}$. Conversely if $f=0$ a.e., every simple $s\le f$ vanishes a.e. and has
integral $0$. (c) $\mu(f=\infty)\le\mu(f\ge t)\le\frac1t\int f$ for all $t$. (d) Let $N$ be null with $f=g$ off $N$; then $\int f=\int f\mathbf 1_{N^c}+\int f\mathbf 1_N=\int f\mathbf 1_{N^c}$ by (b), and the same
for $g$.
\end{proof}

\begin{theorem}[Fatou's lemma]\label{thm:fatou6}
For measurable $f_n\ge0$, $\int\liminf_nf_n\,d\mu\le\liminf_n\int f_n\,d\mu$.
\end{theorem}
\begin{proof}
$g_m=\inf_{n\ge m}f_n$ increases to $\liminf_nf_n$, and $\int g_m\le\int f_n$ for every $n\ge m$, so $\int g_m\le\inf_{n\ge m}\int f_n$. Let $m\to\infty$ and use monotone convergence.
\end{proof}

\begin{definition}
A measurable $f:E\to\R$ (or $\C$) is \emph{integrable} if $\int|f|\,d\mu<\infty$; then $\int f=\int f^+-\int f^-$ for real $f$, and $\int f=\int\operatorname{Re}f+i\int\operatorname{Im}f$ for complex $f$. The integrable functions
form a vector space $\mathcal L^1(\mu)$, on which $f\mapsto\int f$ is linear and $|\int f|\le\int|f|$.
\end{definition}

Linearity follows from Corollary~\ref{cor:linearity6}: if $h=f+g$, then $h^++f^-+g^-=h^-+f^++g^+$, and one integrates both sides. For $|\int f|\le\int|f|$ with complex $f$, choose $|c|=1$
with $c\int f=|\int f|$; then $|\int f|=\int\operatorname{Re}(cf)\le\int|f|$.

\begin{theorem}[dominated convergence]\label{thm:dct6}
Let $f_n$ be measurable with $f_n\to f$ a.e., and $|f_n|\le g$ a.e. for an integrable $g$. Then $f$ is integrable, $\int|f_n-f|\,d\mu\to0$, and $\int f_n\to\int f$.
\end{theorem}
\begin{proof}
Changing everything on a null set (Proposition~\ref{prop:markov}(d)), we may assume convergence and domination everywhere. Then $|f|\le g$, and $2g-|f_n-f|\ge0$. By
Fatou, $\int2g\le\liminf_n\int(2g-|f_n-f|)=\int2g-\limsup_n\int|f_n-f|$. Since $\int g<\infty$, $\limsup_n\int|f_n-f|\le0$. Finally $|\int f_n-\int f|\le\int|f_n-f|$.
\end{proof}

A finite measure space allows constants as dominating functions: \emph{bounded convergence} is the case $g=M$ with $\mu(E)<\infty$. Counting measure on $\N$ turns integrals
into series, and the three theorems into statements about series: $\sum_k\lim_na_{n,k}=\lim_n\sum_ka_{n,k}$ when $|a_{n,k}|\le b_k$ with $\sum_kb_k<\infty$ (Tannery's theorem).

\begin{think}
Let $f_n=n\mathbf 1_{]0,\frac1n[}$ on $[0,1]$ with Lebesgue measure. Then $f_n\to0$ everywhere, and $\int f_n=1$ for all $n$. Which hypothesis of each theorem fails?
\end{think}

Monotone convergence does not apply ($f_n$ is not increasing), Fatou applies and is strict ($0<1$), and dominated convergence fails because $\sup_nf_n$, which is about
$\frac1x$, is not integrable. The mass $1$ escapes into the point $0$, where the limit cannot see it (Figure~\ref{fig:escape}). Section~\ref{sec:ui} makes ``no mass escapes'' precise.

\section{Riemann's integral and Lebesgue's criterion}

Let $f:[a,b]\to\R$ be bounded. For a subdivision $P:a=x_0<x_1<\cdots<x_n=b$, the lower and upper \emph{Darboux sums} are
\[L(f,P)=\sum_i(x_i-x_{i-1})\inf_{[x_{i-1},x_i]}f,\qquad U(f,P)=\sum_i(x_i-x_{i-1})\sup_{[x_{i-1},x_i]}f.\]
Refining $P$ raises $L$ and lowers $U$, and $L(f,P)\le U(f,P')$ for any $P,P'$ (compare both with the common refinement). The function $f$ is \emph{Riemann integrable} if
$\sup_PL(f,P)=\inf_PU(f,P)$, the common value being $\int_a^bf$. This is equivalent to the definition by Riemann sums or by step functions used in Volume~1.

\begin{think}
Is the indicator of $\Q\cap[0,1]$ integrable? Compute its lower and upper Darboux sums for any subdivision, then its Lebesgue integral. Which of the two integrals
can tell that $\Q$ is ``small''?
\end{think}

\begin{theorem}[the integral of a continuous function]\label{thm:riemanncont}
A continuous $f:[a,b]\to\R$ is Riemann integrable. The integral is linear, monotone, additive over adjacent intervals, and $\int_a^bc=c(b-a)$.
\end{theorem}
\begin{proof}
By Heine's theorem (Theorem~\ref{thm:heine}), for $\varepsilon>0$ there is $\delta>0$ with $|f(x)-f(y)|<\varepsilon$ when $|x-y|<\delta$. If all intervals of $P$ are shorter than $\delta$, then on each of them
$\sup f-\inf f\le\varepsilon$, so $U(f,P)-L(f,P)\le\varepsilon(b-a)$. Hence $\sup L=\inf U$. Monotonicity and $\int c=c(b-a)$ are clear from the sums, and additivity from subdivisions that contain the
common endpoint. For linearity, $L(f,P)+L(g,P)\le L(f+g,P)\le U(f+g,P)\le U(f,P)+U(g,P)$, and $L(cf,P)=cL(f,P)$ for $c\ge0$ (with $U$ and $L$ exchanged for $c<0$).
\end{proof}

This is the integral that Volume~1 admitted; everything there (the fundamental theorem of calculus, integration by parts, Taylor's formula) was
derived from the properties listed. Now compare with Lebesgue's integral against $\lambda$ on $[a,b]$.

\begin{theorem}[Lebesgue, 1902]\label{thm:lebesguecriterion}
A bounded function $f:[a,b]\to\R$ is Riemann integrable iff the set of its points of discontinuity is null. Then $f$ is Lebesgue measurable and
$\int_a^bf=\int_{[a,b]}f\,d\lambda$.
\end{theorem}
\begin{proof}
Let $P_n$ be the subdivision of $[a,b]$ into $2^n$ equal intervals, and $\ell_n$, $u_n$ the step functions equal on each open interval of $P_n$ to the infimum and the supremum of $f$ on
its closure (and to $f$ at the points of $P_n$). Then $\ell_n\le f\le u_n$, $\ell_n$ increases and $u_n$ decreases, to limits $\ell\le f\le u$, Borel functions, and since $|f|\le M$, dominated convergence
gives $\int\ell_n\,d\lambda=L(f,P_n)\to\int\ell\,d\lambda$ and $U(f,P_n)\to\int u\,d\lambda$.

\emph{Step 1: $f$ is Riemann integrable iff $U(f,P_n)-L(f,P_n)\to0$.} If so, $\sup L=\inf U$. Conversely, let $P$ have $m$ interior points with $U(f,P)-L(f,P)<\varepsilon$. An interval of
$P_n$ that contains no point of $P$ lies inside an interval of $P$; at most $2m$ intervals of $P_n$ contain a point of $P$, and on each of those $\sup f-\inf f\le2M$. Comparing with the common
refinement $P\cup P_n$, $U(f,P_n)-L(f,P_n)\le U(f,P)-L(f,P)+2m\cdot2M\frac{b-a}{2^n}<2\varepsilon$ for $n$ large.

\emph{Step 2: off the countable set $D=\bigcup_nP_n$, $f$ is continuous at $x$ iff $\ell(x)=u(x)$.} For $x\notin D$, the interval of $P_n$ containing $x$ is a neighbourhood of $x$ shrinking to $x$. If
$f$ is continuous at $x$, $\sup f-\inf f$ on these intervals tends to $0$, so $u(x)=\ell(x)$. If $u(x)-\ell(x)<\varepsilon$, then for $n$ large, $\sup f-\inf f<\varepsilon$ on a neighbourhood of $x$: the
oscillation of $f$ at $x$ is less than $\varepsilon$; so $u(x)=\ell(x)$ means continuity at $x$.

\emph{Conclusion.} By Step~1 and monotone convergence, $f$ is Riemann integrable iff $\int(u-\ell)\,d\lambda=0$, iff $u=\ell$ a.e. (Proposition~\ref{prop:markov}(b)), iff, by Step~2 and since $D$ is countable,
$f$ is continuous a.e. Then $\ell=f=u$ a.e.; $f$ differs from the Borel function $\ell$ on a subset of a null set, so it is Lebesgue measurable, and
$\int_a^bf=\lim L(f,P_n)=\int\ell\,d\lambda=\int f\,d\lambda$.
\end{proof}

So monotone functions, continuous functions, and functions with countably many discontinuities are Riemann integrable. The indicator of
$\Q\cap[0,1]$ is discontinuous everywhere, hence not Riemann integrable, although $\int\mathbf 1_\Q\,d\lambda=0$; Thomae's function (Exercise~\ref{ex:thomae}) is Riemann
integrable with integral $0$. From now on $\int_a^bf\,dx$ denotes the Lebesgue integral; by the theorem, nothing changes for the functions of Volumes~1 to~5.

\emph{Improper integrals.} If $f\ge0$ is continuous on $[a,+\infty[$, then $\int_a^xf\,dt\to\int_{[a,\infty[}f\,d\lambda$ as $x\to\infty$, by monotone convergence: the improper Riemann integral and the Lebesgue integral
agree, finite or not. For signed $f$ they agree when $\int|f|<\infty$, by dominated convergence. But $\int_0^x\frac{\sin t}t\,dt$ has a limit ($\frac\pi2$, Exercise~\ref{ex:dirichlet}) while
$\int_0^\infty\frac{|\sin t|}t\,dt=\infty$ (on $[k\pi,(k+1)\pi]$ the integral is at least $\frac2{(k+1)\pi}$): $\frac{\sin t}t$ is not Lebesgue integrable, and its improper integral is a limit, not an integral.

\section{Modes of convergence}

\begin{definition}
Measurable functions $f_n$ \emph{converge in measure} to $f$ if $\mu(|f_n-f|>\varepsilon)\to0$ for every $\varepsilon>0$, and \emph{in $L^1$} if $\int|f_n-f|\,d\mu\to0$.
\end{definition}

In probability, convergence in measure is convergence in probability, and $L^1$ convergence is $\mathbb E|X_n-X|\to0$.

\begin{proposition}\label{prop:modes}
\begin{enumerate}[label=(\alph*)]
\item Convergence in $L^1$ implies convergence in measure.
\item If $f_n\to f$ in measure, a subsequence converges to $f$ a.e.
\item If $\mu(E)<\infty$, convergence a.e. implies convergence in measure.
\item (Egorov) If $\mu(E)<\infty$ and $f_n\to f$ a.e. (all finite), then for every $\varepsilon>0$ there is $A$ with $\mu(A)<\varepsilon$ such that $f_n\to f$ uniformly on $E\setminus A$.
\end{enumerate}
\end{proposition}
\begin{proof}
(a) Markov: $\mu(|f_n-f|>\varepsilon)\le\frac1\varepsilon\int|f_n-f|$. (b) Choose $n_1<n_2<\cdots$ with $\mu(|f_{n_k}-f|>2^{-k})<2^{-k}$. By Borel--Cantelli, a.e. $x$ has $|f_{n_k}(x)-f(x)|\le2^{-k}$ for all large $k$.
(c) $\mathbf 1_{\{|f_n-f|>\varepsilon\}}\to0$ a.e. and is bounded by $1$; bounded convergence gives $\mu(|f_n-f|>\varepsilon)\to0$. (d) For each $k$, the sets $B_{k,m}=\bigcup_{n\ge m}\{|f_n-f|>\frac1k\}$ decrease in
$m$ to a null set, so by continuity from above some $m_k$ has $\mu(B_{k,m_k})<\varepsilon2^{-k}$. Let $A=\bigcup_kB_{k,m_k}$: $\mu(A)<\varepsilon$, and off $A$, $|f_n-f|\le\frac1k$ for all $n\ge m_k$, for every $k$.
\end{proof}

Egorov's theorem is Littlewood's \emph{third principle}: a.e. convergence is nearly uniform. None of the converses holds: the ``typewriter'' (Exercise~\ref{ex:typewriter})
converges in measure and in $L^1$ but at no point, and $n\mathbf 1_{]0,\frac1n[}$ converges everywhere but not in $L^1$.

\section{Uniform integrability and Vitali's theorem}\label{sec:ui}

\begin{figure}[ht]
\centering
\fig{ch07-escape}
\caption{Left: $f_n=n\mathbf 1_{]0,\frac1n[}$ for $n=1,2,4,8$: each has integral $1$, the limit $0$ has integral $0$; the mass escapes in ever thinner spikes. Right: $g_n=\sqrt n\,\mathbf 1_{]0,\frac1n[}$
also tends to $0$, but $\int_{\{g_n>M\}}g_n\le\frac1{\sqrt n}\le\frac1M$ there: a uniformly integrable family, whose integrals do tend to $0$.}
\label{fig:escape}
\end{figure}

\begin{definition}
Let $\mu$ be finite. A family $\mathcal F$ of integrable functions is \emph{uniformly integrable} if
\[\sup_{f\in\mathcal F}\int_{\{|f|>M\}}|f|\,d\mu\longrightarrow0\qquad\text{as }M\to\infty.\]
\end{definition}

A single integrable $f$ is uniformly integrable, by dominated convergence:
\[
|f|\mathbf 1_{\{|f|>M\}}\to0
\]
almost everywhere, since $f$ is finite almost everywhere, and the truncations are dominated by $|f|$.
Thus any finite family is uniformly integrable, as is any family dominated by one integrable $g$. A uniformly integrable family is bounded in $L^1$: $\int|f|\le M\mu(E)+\int_{\{|f|>M\}}|f|$. In probability ($\mu(E)=1$), the condition reads
$\sup_X\mathbb E[|X|\mathbf 1_{\{|X|>M\}}]\to0$: no mass can hide in the tails.

\begin{lemma}\label{lem:abscont}
If $f$ is integrable, then for every $\varepsilon>0$ there is $\delta>0$ such that $\mu(A)<\delta$ implies $\int_A|f|<\varepsilon$.
\end{lemma}
\begin{proof}
$\int_A|f|\le M\mu(A)+\int_{\{|f|>M\}}|f|$. Choose $M$ so that the second term is less than $\frac\varepsilon2$, then $\delta=\frac\varepsilon{2M}$.
\end{proof}

We can now say exactly when the integrals of a convergent sequence converge. The answer, found by Giuseppe Vitali in 1907 (the
same Vitali who built the non-measurable set), is that no mass may escape: the sequence must converge in measure, and must be
uniformly integrable.

\begin{theorem}[Vitali's convergence theorem]\label{thm:vitali7}
Let $\mu$ be finite, $f_n$ integrable and $f_n\to f$ in measure. Then $f$ is integrable and $\int|f_n-f|\,d\mu\to0$ iff $(f_n)$ is uniformly integrable.
\end{theorem}
\begin{proof}
($\Leftarrow$) Let $C=\sup_n\int|f_n|<\infty$. By Proposition~\ref{prop:modes}(b) a subsequence tends to $f$ a.e., so by Fatou $\int|f|\le C$: $f$ is integrable. Let $\varphi_M(t)=\max(-M,\min(M,t))$,
the truncation at level $M$, which is $1$-Lipschitz, and $|t-\varphi_M(t)|\le|t|\mathbf 1_{\{|t|>M\}}$. Then
\[\begin{aligned}\int|f_n-f|&\le\int|f_n-\varphi_M(f_n)|+\int|\varphi_M(f_n)-\varphi_M(f)|\eqbr+\int|\varphi_M(f)-f|.\end{aligned}\]
The first term is at most $\sup_n\int_{\{|f_n|>M\}}|f_n|$, the third at most $\int_{\{|f|>M\}}|f|$; both are less than $\varepsilon$ for $M$ large. Fix such an $M$. The middle integrand $h_n$ is at most
$2M$ and at most $|f_n-f|$, so $\int h_n\le\varepsilon\mu(E)+2M\mu(|f_n-f|>\varepsilon)$, which is less than $\varepsilon(\mu(E)+1)$ for $n$ large. So $\limsup_n\int|f_n-f|\le\varepsilon(\mu(E)+3)$ for every $\varepsilon$.

($\Rightarrow$) Let $\varepsilon>0$ and $N$ with $\int|f_n-f|<\varepsilon$ for $n\ge N$. Let $\delta$ be given by Lemma~\ref{lem:abscont} for $f$. Since $\sup_n\int|f_n|=C<\infty$, Markov gives $\mu(|f_n|>M)\le\frac CM<\delta$ for
$M$ large, and then for $n\ge N$
\[\int_{\{|f_n|>M\}}|f_n|\le\int|f_n-f|+\int_{\{|f_n|>M\}}|f|<2\varepsilon.\]
The finitely many $f_n$, $n<N$, form a uniformly integrable family. So $\int_{\{|f_n|>M\}}|f_n|<2\varepsilon$ for all $n$ once $M$ is large enough.
\end{proof}

In probability: if $X_n\to X$ in probability, then $\mathbb E|X_n-X|\to0$ iff $(X_n)$ is uniformly integrable, and then $\mathbb E[X_n]\to\mathbb E[X]$. Dominated convergence is the special
case of a dominated family; Exercise~\ref{ex:vallee} gives the practical criterion $\sup_n\mathbb E[|X_n|^p]<\infty$ for some $p>1$. Volume~11 uses the theorem for martingales.

\section{Integrals depending on a parameter}

\begin{theorem}\label{thm:parameter}
Let $T$ be a metric space (for (b), an open interval of $\R$), and $h:E\times T\to\R$ (or $\C$) with $h(\cdot,t)$ integrable for each $t$. Let $H(t)=\int h(x,t)\,d\mu(x)$.
\begin{enumerate}[label=(\alph*)]
\item (Continuity) If $t\mapsto h(x,t)$ is continuous at $t_0$ for a.e. $x$, and $|h(x,t)|\le g(x)$ for all $t$ near $t_0$ with $g$ integrable, then $H$ is continuous at $t_0$.
\item (Differentiation) If $t\mapsto h(x,t)$ is differentiable on $T$ for a.e. $x$, and $|\partial_th(x,t)|\le g(x)$ for all $t\in T$ with $g$ integrable, then $H$ is differentiable on $T$ and
\[H'(t)=\int\partial_th(x,t)\,d\mu(x).\]
\end{enumerate}
\end{theorem}
\begin{proof}
(a) If $t_n\to t_0$, dominated convergence applies to $h(\cdot,t_n)\to h(\cdot,t_0)$. (b) Fix $t$ and $t_n\to t$, $t_n\ne t$. The quotients $q_n(x)=\frac{h(x,t_n)-h(x,t)}{t_n-t}$ tend to $\partial_th(x,t)$ for a.e.
$x$ and, by the mean value theorem (applied to real and imaginary parts), $|q_n(x)|\le2g(x)$. In particular $x\mapsto\partial_th(x,t)$, an a.e.\ limit of the measurable $q_n$, is measurable and bounded by $g$, hence integrable.
Dominated convergence gives $\frac{H(t_n)-H(t)}{t_n-t}=\int q_n\,d\mu\to\int\partial_th(x,t)\,d\mu(x)$, for every sequence $t_n\to t$.
\end{proof}

The hypotheses only ask for a bound uniform in $t$, so in practice one restricts $t$ to a compact subinterval of $T$, where bounds are easier, and
concludes on all of $T$ since differentiability is local. This justifies every application of Feynman's trick in Volume~5 and in the Gaussian and
Gamma integrals of Volume~4. Exercises~\ref{ex:gausscos} and~\ref{ex:dirichlet} are typical.

\begin{beauty}
Lebesgue's idea, ``slice the values, not the variable'', sounds like a minor change of bookkeeping. It made limits commute
with integrals under natural conditions, turned the integral into a tool valid for any measure (lengths, probabilities,
counting, sums of series, surface areas, Haar measures on groups), and made the space of integrable functions complete
(Chapter~10). Not bad for a way of counting coins.
\end{beauty}

\section{Exercises}

\begin{exercise}[$\star$]\label{ex:fatoustrict}
Find measurable $f_n\ge0$ on $(\R,\lambda)$ with (a) $f_n$ decreasing to $0$ but $\int f_n=+\infty$ for all $n$; (b) $f_n\to0$ everywhere, $\int f_n=1$, and all $f_n$ bounded by $1$. Which theorem
fails in each case, and which hypothesis is missing?
\end{exercise}

\begin{exercise}[$\star$]\label{ex:typewriter}
\emph{The typewriter.} Write $n=2^k+j$ with $0\le j<2^k$ and let $f_n=\mathbf 1_{[j2^{-k},(j+1)2^{-k}]}$ on $[0,1]$. Show that $f_n\to0$ in $L^1$ and in measure, but that $(f_n(x))$ converges for no $x$. Find a
subsequence converging a.e.
\end{exercise}

\begin{exercise}[$\star$]\label{ex:riemannalgebra}
Show, with Lebesgue's criterion, that the Riemann integrable functions on $[a,b]$ form an algebra stable under $|\cdot|$, $\max$ and $\min$, and that $g\circ f$ is Riemann integrable
when $f$ is and $g$ is continuous. Is $f\circ g$ always Riemann integrable when $f$ is Riemann integrable and $g$ continuous?
\end{exercise}

\begin{exercise}[$\star\star$]\label{ex:layercake}
\emph{The layer cake formula.} Let $f\ge0$ be measurable. Show that $\int f\,d\mu=\int_0^\infty\mu(f>t)\,dt$. (First for simple $f$, then by monotone convergence.) In probability,
$\mathbb E[X]=\int_0^\infty\mathbb P(X>t)\,dt$ for $X\ge0$, and $\mathbb E[X^p]=\int_0^\infty pt^{p-1}\mathbb P(X>t)\,dt$.
\end{exercise}

\begin{exercise}[$\star\star$]\label{ex:scheffe}
\emph{Scheff\'e's lemma.} Let $f_n,f\ge0$ be integrable with $f_n\to f$ a.e. and $\int f_n\to\int f$. Show that $\int|f_n-f|\to0$. Deduce that if densities $p_n$ of laws on $\R$ converge a.e. to a
density $p$, the laws converge in total variation: $\sup_B|\int_Bp_n-\int_Bp|\to0$.
\end{exercise}

\begin{exercise}[$\star\star$]\label{ex:vallee}
\emph{De la Vall\'ee Poussin's criterion.} Let $\mu$ be finite. Show that a family $\mathcal F$ is uniformly integrable if $\sup_{f\in\mathcal F}\int\Phi(|f|)\,d\mu<\infty$ for some $\Phi:[0,\infty[\to[0,\infty[$ with
$\frac{\Phi(t)}t\to\infty$; for instance $\Phi(t)=t^p$, $p>1$, or $t\log^+t$. (The converse is also true.)
\end{exercise}

\begin{exercise}[$\star\star$]\label{ex:gausscos}
Let $I(t)=\int_0^\infty e^{-x^2}\cos(tx)\,dx$. Justify that $I'(t)=-\int_0^\infty xe^{-x^2}\sin(tx)\,dx$, integrate by parts to get $I'=-\frac t2I$, and deduce $I(t)=\frac{\sqrt\pi}2e^{-t^2/4}$. What does this say about
the characteristic function of a Gaussian (Volume~4)?
\end{exercise}

\begin{exercise}[$\star\star$]\label{ex:rlemma}
\emph{The Riemann--Lebesgue lemma.} Let $f$ be integrable on $\R$. Show that $\int f(x)e^{i\xi x}\,dx\to0$ as $|\xi|\to\infty$. (Approximate $f$ in $L^1$ by a step function, using the regularity of
$\lambda$, Chapter~7.) Volume~1 proved it only for $C^1$ functions.
\end{exercise}

\begin{exercise}[$\star\star\star$]\label{ex:dirichlet}
\emph{The Dirichlet integral.} Let $F(t)=\int_0^\infty e^{-tx}\frac{\sin x}x\,dx$ for $t>0$. (a) Show that $F$ is differentiable on $]0,\infty[$ with $F'(t)=-\frac1{1+t^2}$, and that $F(t)\to0$ as $t\to\infty$; deduce
$F(t)=\frac\pi2-\arctan t$. (b) Show that $\int_0^R\frac{\sin x}x\,dx\to\frac\pi2$ as $R\to\infty$, by proving that $F$ is continuous at $0$ for the improper integral: bound
$\bigl|\int_R^\infty e^{-tx}\frac{\sin x}x\,dx\bigr|$ uniformly in $t\ge0$ by integrating by parts.
\end{exercise}

\section{Hints}

\begin{hint}{ex:fatoustrict}
(a) $\mathbf 1_{[n,\infty[}$. (b) $\mathbf 1_{[n,n+1]}$: mass escapes to infinity rather than into a point.
\end{hint}
\begin{hint}{ex:typewriter}
$\int f_n=2^{-k}$. Every $x$ lies in one interval of each generation and outside the others.
\end{hint}
\begin{hint}{ex:riemannalgebra}
Discontinuities of $f+g$, $fg$, $|f|$, $g\circ f$ lie in the union of those of $f$ and $g$. For $f\circ g$, take $f$ the indicator of $]0,1]$ and $g$ continuous, zero exactly on a fat
Cantor set.
\end{hint}
\begin{hint}{ex:layercake}
For $s=\sum_ia_i\mathbf 1_{A_i}$ with $0<a_1<\cdots<a_m$ and disjoint $A_i$, $\mu(s>t)$ is a step function of $t$. For $s_n\uparrow f$, the sets $\{s_n>t\}$ increase to $\{f>t\}$.
\end{hint}
\begin{hint}{ex:scheffe}
$|f_n-f|=f_n+f-2\min(f_n,f)$, and $\min(f_n,f)\le f$: dominated convergence.
\end{hint}
\begin{hint}{ex:vallee}
On $\{|f|>M\}$, $|f|\le\Phi(|f|)\cdot\sup_{t>M}\frac t{\Phi(t)}$.
\end{hint}
\begin{hint}{ex:gausscos}
The $t$-derivative of $e^{-x^2}\cos tx$ is bounded by $xe^{-x^2}$. For the integration by parts, $-xe^{-x^2}$ is the derivative of $\frac12e^{-x^2}$. $I(0)=\frac{\sqrt\pi}2$ (Volume~1).
\end{hint}
\begin{hint}{ex:rlemma}
For $f=\mathbf 1_{]a,b]}$ compute the integral. Step functions with finitely many steps are dense in $L^1$: approximate $\{f>t\}$ by finite unions of intervals, or use Littlewood's first
principle on the simple functions of Proposition~\ref{prop:measurablefns}(c).
\end{hint}
\begin{hint}{ex:dirichlet}
(a) On $[t_0,\infty[$, dominate by $e^{-t_0x}$, using $|\frac{\sin x}x|\le1$. $\int_0^\infty e^{-tx}\sin x\,dx=\frac1{1+t^2}$. (b) Write down an explicit primitive $\psi_t$ of $e^{-tx}\sin x$ and check that $|\psi_t|\le2$ for all $t\ge0$ and $x\ge0$. Integrate by
parts on $[R,S]$, differentiating $\frac1x$: every term is then $O(\frac1R)$, whatever $t$ is.
\end{hint}

\section{Solutions}

\begin{solution}{ex:fatoustrict}
\idea Monotone convergence downwards needs a finite starting integral; dominated convergence needs a dominating integrable function. Mass can escape to infinity
as well as into a point.
\writeup (a) $f_n=\mathbf 1_{[n,\infty[}$ decreases to $0$ but $\int f_n=\infty$. Monotone convergence holds only upwards; downwards it needs $\int f_1<\infty$ (then it follows from
dominated convergence with $g=f_1$). (b) $f_n=\mathbf 1_{[n,n+1]}$: $f_n\to0$ everywhere, $\int f_n=1$, $0\le f_n\le1$, but $\lambda(\R)=\infty$ so the constant $1$ is not integrable, and $\sup_nf_n=\mathbf 1_{[1,\infty[}$ is not
integrable either. Dominated convergence fails; Fatou holds and is strict.
\end{solution}

\begin{solution}{ex:typewriter}
\idea The bump sweeps $[0,1]$ again and again, getting thinner: small in measure, yet it comes back over every point once per sweep.
\writeup $\int|f_n|=2^{-k}\to0$ as $n\to\infty$ (as $k\to\infty$), so $f_n\to0$ in $L^1$ and in measure (Proposition~\ref{prop:modes}(a)). For each $x\in[0,1]$ and each $k$, $x$ lies in at least one of the
intervals $[j2^{-k},(j+1)2^{-k}]$, so $f_n(x)=1$ for at least one $n$ in each generation $2^k\le n<2^{k+1}$, and $f_n(x)=0$ for at least one other $n$ in that generation once $k\ge1$ (there are $2^k\ge2$
intervals, and $x$ lies in at most two). So $(f_n(x))$ takes both values infinitely often. The subsequence $f_{2^k}=\mathbf 1_{[0,2^{-k}]}$ tends to $0$ at every $x>0$, hence a.e.
\end{solution}

\begin{solution}{ex:riemannalgebra}
\idea Lebesgue's criterion turns integrability into a property of the set of discontinuities, and these sets behave well: a combination is continuous
wherever its ingredients are.
\writeup Let $D_f$ be the set of discontinuities of $f$. If $f,g$ are continuous at $x$, so are $f+g$, $fg$, $|f|$, $\max(f,g)$, $\min(f,g)$, and $h\circ f$ for continuous $h$. So their
discontinuities lie in $D_f\cup D_g$, a null set if $f,g$ are Riemann integrable; they are bounded, and Theorem~\ref{thm:lebesguecriterion} applies. For $f\circ g$: let $K$ be a fat
Cantor set in $[0,1]$ (Exercise~\ref{ex:fatcantor}), $g(x)=\operatorname{dist}(x,K)$, continuous, zero exactly on $K$, and $f=\mathbf 1_{]0,\infty[}$, Riemann integrable on any segment (one
discontinuity). Then $f\circ g=\mathbf 1_{[0,1]\setminus K}$, which is discontinuous at every point of $K$ (since $K$ has empty interior, every point of $K$ is a limit of points outside $K$):
its discontinuities have measure $\lambda(K)=\frac12$, and $f\circ g$ is not Riemann integrable.
\end{solution}

\begin{solution}{ex:layercake}
\idea Integrating $\mu(f>t)$ over $t$ adds up the horizontal slices of the region under the graph; Fubini would say it at once, but monotone convergence is enough.
\writeup Let $s=\sum_{i\le m}a_i\mathbf 1_{A_i}$ with $0=a_0<a_1<\cdots<a_m$ and disjoint $A_i$. For $a_{k-1}\le t<a_k$, $\mu(s>t)=\sum_{i\ge k}\mu(A_i)$, so
\[\int_0^\infty\mu(s>t)\,dt=\sum_{k\le m}(a_k-a_{k-1})\sum_{i\ge k}\mu(A_i)=\sum_i\mu(A_i)\sum_{k\le i}(a_k-a_{k-1})=\sum_ia_i\mu(A_i).\]
For $f\ge0$, take simple $s_n\uparrow f$ (Proposition~\ref{prop:measurablefns}(c)). For each $t$, $\{s_n>t\}$ increases to $\{f>t\}$ (if $f(x)>t$, then $s_n(x)>t$ for large $n$), so $\mu(s_n>t)\uparrow\mu(f>t)$. Two
applications of monotone convergence, in $t$ for $\lambda$ and in $x$ for $\mu$, give $\int f\,d\mu=\lim\int s_n\,d\mu=\lim\int_0^\infty\mu(s_n>t)\,dt=\int_0^\infty\mu(f>t)\,dt$. (The functions $t\mapsto\mu(s_n>t)$ are
non-increasing, hence Borel.) For $\mathbb E[X^p]$, apply this to $X^p$ and substitute $t=u^p$ ($\{X^p>u^p\}=\{X>u\}$; the change of variables for monotone $C^1$ maps on intervals holds for
non-negative measurable integrands by approximation, or see Chapter~9).
\end{solution}

\begin{solution}{ex:scheffe}
\idea If the integrals converge and the functions converge pointwise, the mass of $f_n$ cannot escape: what $f_n$ has above $f$ is compensated by what it lacks below, and
the part below is dominated by $f$.
\writeup $|f_n-f|=f_n+f-2\min(f_n,f)$. Since $0\le\min(f_n,f)\le f$ and $\min(f_n,f)\to f$ a.e., dominated convergence gives $\int\min(f_n,f)\to\int f$. Hence $\int|f_n-f|\to\int f+\int f-2\int f=0$. For
densities, $|\int_Bp_n-\int_Bp|\le\int|p_n-p|\to0$ uniformly in $B$.
\end{solution}

\begin{solution}{ex:vallee}
\idea Superlinear growth of $\Phi$ means that large values of $|f|$ cost much more in $\int\Phi(|f|)$ than in $\int|f|$; a bound on the first leaves little room for tails.
\writeup Let $C=\sup_f\int\Phi(|f|)$ and $\eta(M)=\sup_{t>M}\frac t{\Phi(t)}$, which tends to $0$ since $\frac{\Phi(t)}t\to\infty$ (and $\Phi(t)>0$ for $t$ large). On $\{|f|>M\}$, $|f|\le\eta(M)\Phi(|f|)$, so
$\int_{\{|f|>M\}}|f|\le\eta(M)C\to0$ uniformly in $f$. For $\Phi(t)=t^p$, $\eta(M)=M^{1-p}$; for $t\log^+t$, $\eta(M)=\frac1{\log M}$.
\end{solution}

\begin{solution}{ex:gausscos}
\idea Differentiate under the integral sign, which needs one integrable bound for the derivative, uniform in $t$; then an integration by parts turns the
derivative back into $I$ itself.
\writeup The function $h(x,t)=e^{-x^2}\cos(tx)$ satisfies \[|\partial_th(x,t)|=|xe^{-x^2}\sin(tx)|\le xe^{-x^2},\] which is integrable on $[0,\infty[$, so Theorem~\ref{thm:parameter}(b) applies on all of $\R$. Integrating by parts on $[0,R]$
with $u=\sin(tx)$ and $v'=-xe^{-x^2}$, $v=\frac12e^{-x^2}$:
\[\begin{aligned}I'(t)&=\Bigl[\tfrac12e^{-x^2}\sin(tx)\Bigr]_0^\infty-\frac t2\int_0^\infty e^{-x^2}\cos(tx)\,dx\eqbr=-\frac t2I(t).\end{aligned}\]
So $(e^{t^2/4}I)'=0$, and $I(t)=I(0)e^{-t^2/4}=\frac{\sqrt\pi}2e^{-t^2/4}$. Since the integrand is even in $x$, $\int_\R e^{-x^2}e^{itx}\,dx=\sqrt\pi e^{-t^2/4}$; substituting $x=\frac y{\sqrt2}$, the standard
Gaussian law has characteristic function $\mathbb E[e^{itZ}]=e^{-t^2/2}$, as Volume~4 found by another route.
\end{solution}

\begin{solution}{ex:rlemma}
\idea Oscillation kills the integral of an indicator of an interval at speed $\frac1{|\xi|}$; step functions are dense in $L^1$, and the map $f\mapsto\int fe^{i\xi x}$ is $1$-Lipschitz for the $L^1$
norm, uniformly in $\xi$.
\writeup For $f=\mathbf 1_{]a,b]}$, $|\int_a^be^{i\xi x}\,dx|=|\frac{e^{i\xi b}-e^{i\xi a}}{i\xi}|\le\frac2{|\xi|}\to0$; by linearity the same holds for step functions (finite combinations of such indicators).
Now let $f\in\mathcal L^1$ and $\varepsilon>0$. Splitting $f$ into real and imaginary, positive and negative parts, we may take $f\ge0$. By Proposition~\ref{prop:measurablefns}(c) and monotone
convergence, a simple function $s=\sum_ia_i\mathbf 1_{A_i}\le f$ has $\int(f-s)<\varepsilon$, and the $A_i$ have finite measure. By Littlewood's first principle (after Proposition~\ref{prop:regularity}) each $A_i$
is a finite union of intervals up to a set of measure $\frac\varepsilon{m|a_i|}$, so a step function $\phi$ has $\int|s-\phi|<\varepsilon$. Then $\int|f-\phi|<2\varepsilon$, and
\[\Bigl|\int fe^{i\xi x}\Bigr|\le\int|f-\phi|+\Bigl|\int\phi e^{i\xi x}\Bigr|<2\varepsilon+\Bigl|\int\phi e^{i\xi x}\Bigr|,\]
whose last term tends to $0$. So \[\limsup_{|\xi|\to\infty}\Bigl|\int fe^{i\xi x}\Bigr|\le2\varepsilon\] for every $\varepsilon>0$, and the limit is $0$.
\end{solution}

\begin{solution}{ex:dirichlet}
\idea The factor $e^{-tx}$ makes everything absolutely integrable for $t>0$, so Feynman's trick applies there and gives $F$ exactly. The real work is at $t=0$, where
$\frac{\sin x}x$ is not integrable: an integration by parts shows that the tails are small uniformly in $t$, which lets $t\to0$.
\writeup (a) For $t\ge t_0>0$, $|\partial_t(e^{-tx}\frac{\sin x}x)|=|e^{-tx}\sin x|\le e^{-t_0x}$, integrable, so $F$ is differentiable on $]t_0,\infty[$ for each $t_0$, with
\[F'(t)=-\int_0^\infty e^{-tx}\sin x\,dx=-\operatorname{Im}\int_0^\infty e^{(-t+i)x}\,dx=-\operatorname{Im}\frac1{t-i}=-\frac1{1+t^2}.\]
Also $|F(t)|\le\int_0^\infty e^{-tx}\,dx=\frac1t\to0$. Since $\frac\pi2-\arctan t$ has the same derivative and also tends to $0$, $F(t)=\frac\pi2-\arctan t$.
(b) Let $G_R(t)=\int_0^Re^{-tx}\frac{\sin x}x\,dx$ for $t\ge0$, a proper integral, continuous in $t$ by Theorem~\ref{thm:parameter}(a) (dominated by $1$ on $[0,R]$). Let $\psi_t(x)=-e^{-tx}\frac{t\sin x+\cos x}{1+t^2}$, a
primitive of $e^{-tx}\sin x$, with $|\psi_t(x)|\le\frac{t+1}{1+t^2}\le2$. Integrating by parts on $[R,S]$,
\[\int_R^Se^{-tx}\frac{\sin x}x\,dx=\Bigl[\frac{\psi_t(x)}x\Bigr]_R^S+\int_R^S\frac{\psi_t(x)}{x^2}\,dx,\]
so $|\int_R^Se^{-tx}\frac{\sin x}x\,dx|\le\frac2S+\frac2R+\frac2R\le\frac6R$ for all $S>R$ and all $t\ge0$. Hence the improper integrals $\int_0^\infty e^{-tx}\frac{\sin x}x\,dx$ exist for $t\ge0$, equal $F(t)$ for $t>0$, and
$|F(t)-G_R(t)|\le\frac6R$ uniformly in $t\ge0$ (with $F(0)$ the improper integral). A uniform limit of continuous functions is continuous, so $F$ is continuous at $0$:
$\int_0^\infty\frac{\sin x}x\,dx=F(0)=\lim_{t\to0^+}(\frac\pi2-\arctan t)=\frac\pi2$.
\end{solution}

\section*{Sources and further reading}

\begin{itemize}
\item H.~Lebesgue, \emph{Le\c cons sur l'int\'egration et la recherche des fonctions primitives}, Gauthier-Villars, 1904. The integral, the criterion
for Riemann integrability, and the comparison with sorting coins appears in his later lecture \emph{Sur le d\'eveloppement de la notion
d'int\'egrale} (1926).
\item G.~B.~Folland, \emph{Real Analysis}, 2nd ed., Wiley, 1999, Chapter~2. Integration, modes of convergence, Egorov.
\item W.~Rudin, \emph{Real and Complex Analysis}, 3rd ed., McGraw--Hill, 1987, Chapter~1.
\item P.~Billingsley, \emph{Probability and Measure}, 3rd ed., Wiley, 1995, Sections~15--16. Uniform integrability and Vitali's theorem in
probabilistic language.
\item D.~Williams, \emph{Probability with Martingales}, Cambridge University Press, 1991, Chapters~5--6 and~13. Monotone convergence, Fatou,
dominated convergence and uniform integrability, with a probabilist's eye.
\item G.~Vitali, \emph{Sull'integrazione per serie}, Rendiconti del Circolo Matematico di Palermo 23 (1907), 137--155.
\item \emph{For the stories in the introduction.} The paying-a-debt story is paraphrased above, for the reason given there. B.~Riemann, \emph{\"Uber die Darstellbarkeit einer Function durch
eine trigonometrische Reihe}, Habilitationsschrift, G\"ottingen, 1854; B.~Levi, ``Sopra l'integrazione delle serie'', \emph{Rend.\ Ist.\ Lombardo} 39 (1906), 775--780; P.~Fatou,
``S\'eries trigonom\'etriques et s\'eries de Taylor'', \emph{Acta Math.} 30 (1906), 335--400. For the whole history: T.~Hawkins, \emph{Lebesgue's Theory of Integration: Its
Origins and Development}, 2nd ed., Chelsea, 1975.
\end{itemize}
